\documentclass[10pt]{article}
\usepackage[T1]{fontenc}
\usepackage[utf8]{inputenc}
\usepackage[french]{babel}
\usepackage{amsmath,amssymb,mathrsfs,mathpazo}
\usepackage{geometry}
\geometry{a4paper,margin=2cm}
\usepackage{graphicx,array,tabularx,tabularray,colortbl,multirow,multicol,booktabs,enumitem,pifont,ulem,fancybox,diagbox,cancel}
\usepackage{pstricks,pstricks-add,pst-plot,pst-tree,pst-node,pst-text,pst-func,pst-math,pst-bspline,pst-3dplot,pst-eucl,pst-solides3d}
\let\pscalculate\undefined
\usepackage{pst-bezier}
\usepackage{tikz,pgfplots,tkz-tab}
\usetikzlibrary{arrows,arrows.meta,patterns,calc,positioning,trees,decorations.markings,shapes.geometric,shapes.misc,shapes.arrows,intersections,decorations.pathmorphing,decorations.pathreplacing,backgrounds,shadows,babel}
\pgfplotsset{compat=1.18}
\usepgfplotslibrary{fillbetween}
\usepackage[np]{numprint}
\usepackage{listings,makecell,pas-tableur,twoopt}
\let\barre\undefined
\usepackage[pstricks]{tablvar}
\definecolor{darkgreen}{RGB}{0,100,0}
\providecommand{\red}{\color{red}}
\providecommand{\blue}{\color{blue}}
\providecommand{\green}{\color{green}}
\providecommand{\black}{\color{black}}
\providecommand{\white}{\color{white}}
\providecommand{\gray}{\color{gray}}
\providecommand{\cyan}{\color{cyan}}
\providecommand{\magenta}{\color{magenta}}
\providecommand{\yellow}{\color{yellow}}
\providecommand{\violet}{\color{violet}}
\providecommand{\orange}{\color{orange}}
\providecommand{\euro}{\texteuro}
\providecommand{\EUR}{\texteuro}
\providecommand{\ieme}{\textsuperscript{e}}
\providecommand{\degres}{\ensuremath{^\circ}}
\providecommand{\degree}{\ensuremath{^\circ}}
\providecommand{\vv}[1]{\overrightarrow{#1}}
\providecommand{\Coord}[1]{#1}
\providecommand{\decosix}{}
\providecommand{\starredbullet}{\ensuremath{\bullet}}
\pagestyle{empty}
\setlength{\parindent}{0pt}
\usepackage{adjustbox}\DeclareUnicodeCharacter{2236}{:}\DeclareUnicodeCharacter{2190}{\ensuremath{\leftarrow}}
\providecommand{\R}{\mathbb{R}}
\providecommand{\N}{\mathbb{N}}
\providecommand{\Z}{\mathbb{Z}}
\providecommand{\Q}{\mathbb{Q}}
\providecommand{\D}{\mathbb{D}}
\providecommand{\C}{\mathbb{C}}
\providecommand{\e}{{\mathrm{e}}}
\providecommand{\vect}[1]{\overrightarrow{#1}}
\providecommand{\vectt}[1]{\overrightarrow{\mathrm{#1}}}
\providecommand{\vvt}[1]{\overrightarrow{\mathrm{#1}}}
\providecommand{\np}[2][]{{#2~#1}}
\providecommand{\vv}[1]{{\overrightarrow{#1}}}
\renewcommand{\d}{{\mathrm{d}}}
\providecommand{\degres}{{^\circ}}
\providecommand{\degre}{{^\circ}}
\providecommand{\blue}{\color{blue}}
\providecommand{\red}{\color{red}}
\providecommand{\enskip}{\quad}
\providecommand{\llbracket}{\lbrack\!\lbrack}
\providecommand{\rrbracket}{\rbrack\!\rbrack}
\providecommand{\euro}{{\text{€}}}
\providecommand{\no}{{\text{n°}}}
\providecommand{\interval}[2]{{[#1;#2]}}
\providecommand{\intervFO}[2]{{[#1;#2[}}
\providecommand{\dotfill}{\ldots}

\providecommand{\ds}{\displaystyle}
\providecommand{\textding}[1]{#1}
\providecommand{\Oij}{\(\mathrm{O};\vec{\imath},\vec{\jmath}\)}
\providecommand{\Oijk}{\(\mathrm{O};\vec{\imath},\vec{\jmath},\vec{k}\)}
\providecommand{\Ouv}{\(\mathrm{O};\vec{u},\vec{v}\)}
\providecommand{\intervFF}[2]{\ensuremath{[#1;#2]}}
\renewcommand{\euro}{\ifmmode\text{\texteuro}\else\texteuro\fi}
\DeclareUnicodeCharacter{00B0}{\ifmmode^{\circ}\else\textdegree\fi}
\DeclareUnicodeCharacter{20AC}{\ifmmode\text{\texteuro}\else\texteuro\fi}
% Énoncés : Éducation nationale. Transcription et corrigés : APMEP, Denis Vergès et contributeurs. Découpage et classement : CoopMaths, Sébastien Lozano, Éric Elter et contributeurs. Adaptation de la présentation : Maths Explorer.
% CC BY-SA : https://coopmaths.fr/www/about/
% https://coopmaths.fr/alea/static/eam/2026/tex/eam_specifique_2026_09_metropole_automatismes.tex
\begin{document}
\noindent\textit{APMEP / CoopMaths. CC BY-SA. Présentation adaptée par Maths Explorer.}\par\medskip
\textbf{Pour cette première partie, aucune justification n'est demandée et une seule réponse est possible par question. Pour chaque question, reporter son numéro sur la copie et indiquer la réponse.}

Une réponse fausse ou l'absence de réponse n'enlève aucun point.

\subsection*{Question 1}

Un élève a obtenu les notes suivantes : 11 ; 12 ; 15 ; 9 ; 17 ; 9 ; 8.

La médiane de ses notes est :

\medskip
\begin{tblr}{colspec = {@{}*{4}{X}}}
	\textbf{A.}\quad 9 &
	\textbf{B.}\quad 11 &
	\textbf{C.}\quad 81 &
	\textbf{D.}\quad 12 \\
\end{tblr}

\subsection*{Question 2}

\begin{minipage}[t]{0.6\linewidth}
	La fonction affine $f$ représentée ci-contre a pour expression :

	\medskip
	\begin{tblr}{colspec = {@{}X}}
		\textbf{A.}\quad $f(x) = 2x+6$ \\
		\textbf{B.}\quad $f(x) = -2x+6$ \\
		\textbf{C.}\quad $f(x) = 6x+2$ \\
		\textbf{D.}\quad $f(x) = -3x+6$ \\
	\end{tblr}
\end{minipage}
\hfill
\begin{tikzpicture}[x=0.75cm,y=0.75cm,>=stealth,baseline={(current bounding box.north)}]
	\draw[step=1,gray!75,line width=0.5pt] (-1.5,-1.5) grid (3.5,6.5);
	\draw[->,line width=1pt] (-1.5,0) -- (3.5,0);
	\draw[->,line width=1pt] (0,-1.5) -- (0,6.5);
	\foreach \x in {-1,1,2} \draw (\x,2pt) -- (\x,-2pt) node[below] {\footnotesize \np{\x}};
	\foreach \y in {-1,1,2,3,4,5,6} \draw (2pt,\y) -- (-2pt,\y) node[left] {\footnotesize \np{\y}};
	\draw (-2pt,-2pt) node[below left] {\footnotesize 0};
	\clip (-1.5,-1.5) rectangle (3.5,6.5);
	\draw[line width=1.2pt,red] plot[domain=-0.5:2.5] (\x,{-3*\x+6});
\end{tikzpicture}

\subsection*{Question 3}

L'expression $x^2 - 9$ correspond au développement de :

\medskip
\begin{tblr}{colspec = {@{}*{4}{X}}}
	\textbf{A.}\quad $(x-3)^2$ &
	\textbf{B.}\quad $(x-3)(x+3)$ &
	\textbf{C.}\quad $9(x-1)^2$ &
	\textbf{D.}\quad $(x-9)(x+1)$ \\
\end{tblr}

\subsection*{Question 4}

30\,\% de 12 est égal à :

\medskip
\begin{tblr}{colspec = {@{}*{4}{X}}}
	\textbf{A.}\quad 3,6 &
	\textbf{B.}\quad 36 &
	\textbf{C.}\quad 4 &
	\textbf{D.}\quad 6 \\
\end{tblr}

\subsection*{Question 5}

L'expression $A = \dfrac{\frac{3}{7}}{\frac{4}{5}}$ est égale à :

\medskip
\begin{tblr}{colspec = {@{}*{4}{X}}}
	\textbf{A.}\quad $\dfrac{12}{35}$ &
	\textbf{B.}\quad $\dfrac{15}{32}$ &
	\textbf{C.}\quad $\dfrac{15}{28}$ &
	\textbf{D.}\quad $\dfrac{28}{15}$ \\
\end{tblr}

\subsection*{Question 6}

Sachant que 1 mètre correspond à environ 1,1 yard, 300 mètres représentent environ :

\medskip
\begin{tblr}{colspec = {@{}*{4}{X}}}
	\textbf{A.}\quad 270 yards &
	\textbf{B.}\quad 303 yards &
	\textbf{C.}\quad 330 yards &
	\textbf{D.}\quad 333 yards \\
\end{tblr}

\subsection*{Question 7}

Le nombre $2 \times 2^5$ est égal à :

\medskip
\begin{tblr}{colspec = {@{}*{4}{X}}}
	\textbf{A.}\quad $4^5$ &
	\textbf{B.}\quad $2^{10}$ &
	\textbf{C.}\quad $4^{10}$ &
	\textbf{D.}\quad $2^6$ \\
\end{tblr}

\subsection*{Question 8}

On note $T_F$ la température en degrés Fahrenheit et $T_C$ la température en degrés Celsius.

On a la relation : $T_F = 1,8\,T_C + 32$.

Si la température $T_C$ est de $-10$~\textdegree C, quelle est la température en degrés Fahrenheit ?

\medskip
\begin{tblr}{colspec = {@{}*{4}{X}}}
	\textbf{A.}\quad 30,2 &
	\textbf{B.}\quad 14 &
	\textbf{C.}\quad 50 &
	\textbf{D.}\quad $-50$ \\
\end{tblr}

\end{document}
